\providecommand{\Res}[2]{\qty[round-mode=figures,round-precision=3]{\fpeval{#1}}{#2}}
\providecommand{\Val}[1]{\num[round-mode=figures,round-precision=4]{\fpeval{#1}}}
\providecommand{\ResSci}[2]{\qty[round-mode=figures,round-precision=3,exponent-mode=scientific]{\fpeval{#1}}{#2}}
\subsubsection*{Question a) Diagramme des forces}
Deux forces : le poids $mg$ vers le bas et la force électrique $|q|E$ vers le haut. Pour une charge négative, la force est opposée au champ : le champ doit donc pointer \textbf{vers le bas}.

\subsubsection*{Question b) Charge de la gouttelette}
\begin{enumerate}[label=\arabic*)]
    \item \textbf{Variables connues et recherchées :} $m = \qty{\v{m}e-15}{kg}$, $g = \qty{9.81}{m/s^2}$, $E = \qty[round-mode=figures,round-precision=4]{\fpeval{\v{m}*1e-15*9.81/(\v{n}*1.602e-19)}}{N/C}$ ; variable recherchée : $|q|$
    \item \textbf{Équations pertinentes :} à l'équilibre :
    \[ |q|E = mg \]
    \item \textbf{Isolation de la variable recherchée :}
    \[ |q| = \frac{mg}{E} \]
    \item \textbf{Résolution numérique :}
    \[ |q| = \frac{\num{\v{m}e-15} \cdot \num{9.81}}{\num[round-mode=figures,round-precision=4]{\fpeval{\v{m}*1e-15*9.81/(\v{n}*1.602e-19)}}} = \ResSci{\v{n}*1.602e-19}{C} \]
    \item \textbf{Réponse :} $q = \mathbf{-\ResSci{\v{n}*1.602e-19}{C}}$.
\end{enumerate}

\subsubsection*{Question c) Nombre d'électrons}
\[ n = \frac{|q|}{e} = \frac{\ResSci{\v{n}*1.602e-19}{}}{\num{1.602e-19}} = \num{\v{n}} \]
La gouttelette porte \textbf{\num{\v{n}} électrons} en excès (un nombre entier, ce qui montre que la charge est quantifiée).
